\begin{theorem}[Endpoint weak-type $(1,1)$ bound for the Gaussian vector Riesz transform]
Let $\mathcal{M}(\gamma_n;\mathbb{R}^n)$ denote the space of
$\gamma_n$-almost-everywhere equivalence classes of measurable
$\mathbb{R}^n$-valued functions on $\mathbb{R}^n$.  There is exactly one linear
map
\[
  T : L^1(\gamma_n;\mathbb{R}) \longrightarrow \mathcal{M}(\gamma_n;\mathbb{R}^n)
\]
with the following two properties:
\begin{enumerate}
\item $T$ extends the transform of \noderef{RieszTransform} from $L^2$, i.e.
$T(\iota v) = R_\gamma v$ for every $v \in L^2(\gamma_n)$, with $\iota$ the
inclusion of \noderef{ToL1}; and
\item $T$ satisfies the weak-type bound
\[
  \sup_{t>0}\; t\,\gamma_n\bigl(\{x : |Tf(x)|_{\mathbb{R}^n} > t\}\bigr)
    \;\leq\; 2\,\|f\|_{L^1(\gamma_n)}
  \qquad\text{for every } f \in L^1(\gamma_n;\mathbb{R}) .
\]
\end{enumerate}
The constant $2$ does not depend on $n$.  Property (2) is precisely boundedness
from $L^1(\gamma_n)$ into $L^{1,\infty}(\gamma_n;\mathbb{R}^n)$ with quasi-norm at
most $2$, hence continuity of $T$ for the weak-$L^1$ quasi-norm; so the statement
is the assertion that $R_\gamma$ has a unique continuous linear extension from
$L^1(\gamma_n)$ to $L^{1,\infty}(\gamma_n;\mathbb{R}^n)$ satisfying the bound with
constant $2$, and uniqueness is asserted in the larger class of all linear maps
into $\mathcal{M}(\gamma_n;\mathbb{R}^n)$ satisfying (1) and (2).
\end{theorem}
\begin{proof}
Existence is \noderef{L1Extension}, which supplies a linear
$T : L^1(\gamma_n) \to \mathcal{M}(\gamma_n;\mathbb{R}^n)$ satisfying (1) and (2).

For uniqueness, let $T$ and $T'$ both be linear and satisfy (1) and (2), and fix
$f \in L^1(\gamma_n)$.  Choose $f_k \in L^2(\gamma_n)$ with $f_k \to f$ in
$L^1(\gamma_n)$: for instance the truncations
$f_k = \max\{-k,\min\{f,k\}\}$, which are bounded, hence in $L^2(\gamma_n)$ since
$\gamma_n$ is a probability measure (\noderef{GaussIsProb}), and converge to $f$
in $L^1$ by dominated convergence.  (That $L^2(\gamma_n)$ is dense in
$L^1(\gamma_n)$ also follows from \noderef{L2IntoL1} together with
\noderef{HermiteL2Complete}; the explicit truncations suffice here.)

By (1), $T(\iota f_k) = R_\gamma f_k = T'(\iota f_k)$ for every $k$, so
\[
  T f - T' f = T(f - \iota f_k) - T'(f - \iota f_k)
\]
by linearity of $T$ and $T'$.  Let $\varepsilon > 0$.  If
$|Tf(x) - T'f(x)| > 2\varepsilon$ then at least one of
$|T(f-\iota f_k)(x)| > \varepsilon$ and $|T'(f-\iota f_k)(x)| > \varepsilon$ holds,
so by subadditivity and (2) applied to each of $T$ and $T'$ with the input
$f - \iota f_k$ and the level $\varepsilon$,
\[
  \gamma_n\bigl(\{|Tf - T'f| > 2\varepsilon\}\bigr)
    \leq \frac{2}{\varepsilon}\|f - \iota f_k\|_{L^1(\gamma_n)}
       + \frac{2}{\varepsilon}\|f - \iota f_k\|_{L^1(\gamma_n)}
    = \frac{4}{\varepsilon}\|f - \iota f_k\|_{L^1(\gamma_n)} .
\]
Letting $k \to \infty$, the right-hand side tends to $0$, so
$\gamma_n(\{|Tf-T'f| > 2\varepsilon\}) = 0$ for every $\varepsilon > 0$.  Taking
$\varepsilon = 1/j$ and using that
$\{Tf \neq T'f\} = \bigcup_{j\geq1}\{|Tf-T'f| > 2/j\}$ is a countable union of
null sets, we get $Tf = T'f$ $\gamma_n$-almost everywhere, i.e.\ $Tf = T'f$ in
$\mathcal{M}(\gamma_n;\mathbb{R}^n)$.  As $f \in L^1(\gamma_n)$ was arbitrary,
$T = T'$.
\end{proof}
